% File: Volume-I-Foundations/chapters/ch03-shifted-symplectic.tex

\chapter{Shifted Symplectic Structures}
\label{chap:shifted-symplectic}

This chapter develops the basic notions of shifted symplectic geometry used to
formulate the moduli stack $\Plenum$ of lamphrodynamic configurations as a
$(-1)$--shifted symplectic derived Artin stack. We assume familiarity with the
cotangent complex of a derived stack and its dual tangent complex as introduced
in \cref{chap:inf-cat}. Throughout, we follow \cite{PTVV2013}, \cite{Calaque2015}
and related sources.

%--------------------------------------------------------------------
\section{Preliminaries: Classical Symplectic Geometry}
%--------------------------------------------------------------------

We begin with a brief review of classical symplectic geometry to fix notation
and highlight the generalization to derived settings.

\begin{definition}
A \emph{symplectic manifold} is a pair $(X,\omega)$ where $X$ is a smooth
manifold and $\omega \in \Omega^2(X)$ is a closed (i.e.\ $\dd\omega = 0$) and
nondegenerate $2$--form.
\end{definition}

A central example is $T^{\ast}Q$ for a smooth manifold $Q$, equipped with the
canonical $2$--form $\omega = \dd p_i \wedge \dd q^i$. Classical Hamiltonian
mechanics is framed on $(X,\omega)$ with Hamiltonian $H\in C^\infty(X)$.

In gauge theory and in field theory, however, the moduli space of classical
solutions is typically singular, infinite dimensional, and equipped with
higher automorphisms, so classical manifolds are no longer adequate. A
homotopical refinement is needed, leading to the notion of shifted symplectic
structures.

%--------------------------------------------------------------------
\section{Graded Geometry and Cohomological Degree}
%--------------------------------------------------------------------

In derived geometry, differential forms are replaced by cohomological complexes
and their degrees must be tracked carefully. A $k$--shifted symplectic form is,
informally, a closed $2$--form of cohomological degree $k$ on the derived
stack. Heuristically:
\[
  \omega \in H^{k}\big(\mathcal{A}^{2,cl}(F)\big),
\]
where $F$ is a derived Artin stack and $\mathcal{A}^{2,cl}(F)$ is the sheaf of
closed $2$--forms on $F$.

\begin{definition}
Let $F$ be a derived Artin stack. A \emph{$k$--shifted symplectic structure} on
$F$ is a closed $2$--form $\omega$ of cohomological degree $k$, represented by
a morphism of complexes
\[
  \omega: \wedge^2 L_F \longrightarrow \mathcal{O}_F[k],
\]
where $L_F$ is the cotangent complex of $F$, such that the induced pairing
$L_F^\vee \otimes L_F^\vee \to \mathcal{O}_F[k]$ is nondegenerate in the
derived sense.
\end{definition}

Here $L_F^\vee = T_F$ generalizes the tangent bundle. The condition that $L_F$
be perfect (locally of finite Tor dimension) ensures that $T_F$ behaves well as
a dual.

\Intuition
The shift $k$ controls the ``oddness'' or ``evenness'' of the symplectic form.
Classical symplectic structures correspond to $k=0$. The case $k=-1$
corresponds to odd symplectic geometry and plays the role of BV phase spaces.

%--------------------------------------------------------------------
\section{$(-1)$--Shifted Symplectic Structures and BV}
%--------------------------------------------------------------------

The case $k=-1$ is the setting of the Batalin–Vilkovisky (BV) formalism for
gauge theories. In this context, the derived stack $F$ encodes classical fields
and gauge symmetries; the shifted symplectic form $\omega$ is odd; and the BV
bracket
\[
  \{-,-\}: \mathcal{O}_F \otimes \mathcal{O}_F \to \mathcal{O}_F[-(-1)] = \mathcal{O}_F[1]
\]
is defined from $\omega$.

\begin{definition}
Given a $(-1)$--shifted symplectic structure $\omega$ on $F$, the
\emph{Batalin–Vilkovisky bracket} $\{-,-\}$ on functions (or, more precisely,
on appropriate cochains) is the unique odd Poisson bracket induced by $\omega$.
\end{definition}

The BV mechanism encodes gauge symmetries as degenerate directions of the
action, formally resolved by ghosts and higher homotopies. The \emph{classical
master equation}
\[
  \{S,S\} = 0
\]
then expresses invariance under these symmetries. Later, in
\cref{chap:aksz-bv,chap:aksz-rsvp}, we will show how the lamphrodynamic action
$S_{\rsvpabbr}$ arises from an AKSZ construction on the moduli stack
$\Plenum$ equipped with a $(-1)$--shifted symplectic structure.

%--------------------------------------------------------------------
\section{Closedness and Nondegeneracy in the Derived Sense}
%--------------------------------------------------------------------

Closedness of $\omega$ is expressed in terms of the differential and the
cohomological structure on $L_F$. Nondegeneracy means that the induced morphism
\[
  T_F \longrightarrow L_F[-k]
\]
is an isomorphism in the derived category. For $k=-1$, this becomes
\[
  T_F \longrightarrow L_F[1],
\]
which is automatically compatible with the BV grade.

\begin{remark}
For $k<0$, the shifted symplectic form is ``odd'' and induces higher order
compatibilities among tangent and cotangent complexes, matching precisely the
structure required for the BV bracket.
\end{remark}

Closedness means that $\omega$ lifts to a cocycle in a suitable truncated
de~Rham complex on $F$. We do not reproduce the full homotopical definition here
(see \cite{PTVV2013}), but emphasize that closedness is a global condition
involving sheaf descent and higher coherences.

%--------------------------------------------------------------------
\section{Examples and Fundamental Results}
%--------------------------------------------------------------------

Key results from Pantev–Toën–Vaquié–Vezzosi \cite{PTVV2013} include:

\begin{theorem}[PTVV]
Let $X$ be a smooth scheme and let $\Tstar{k}X$ be the $k$--shifted cotangent
stack. Then $\Tstar{k}X$ carries a canonical $k$--shifted symplectic structure.
\end{theorem}

This theorem provides the fundamental example of shifted symplectic structures
and plays a role analogous to the cotangent bundle in classical mechanics. In
derived geometry, one generalizes the phase space to $\Tstar{k}X$ for any
integer $k$, with $k=-1$ corresponding to BV theory.

Another key result concerns mapping stacks:

\begin{theorem}[PTVV]
Let $M$ be an oriented smooth manifold of dimension $d$, and let $F$ be a
derived Artin stack with an $n$--shifted symplectic structure. Then the
mapping stack $\Map(M,F)$ carries an $(n-d)$--shifted symplectic structure.
\end{theorem}

This result is central to the AKSZ construction, where $M$ is typically the
source manifold (e.g.\ the spacetime worldvolume) and $F$ is the target stack
encoding the fields. For \rsvpabbr{}, $M$ will often be spacetime or a
spacelike hypersurface, and $F$ the lamphrodynamic stack $\Plenum$.

\Intuition
Shifted symplecticity propagates from target to mapping stack, reducing the
shift by the dimension of the domain. For a $3$--dimensional worldvolume and
an $n$--shifted target, the resulting structure is $(n-3)$--shifted.

%--------------------------------------------------------------------
\section{Shifted Symplectic Forms from Field Content}
%--------------------------------------------------------------------

In \rsvpabbr{}, the $(-1)$--shifted symplectic structure arises from the
lamphrodynamic triplet $(\PhiField,\vField,\SField)$ and their coupling to
geometry. Roughly speaking, the moduli stack $\Plenum$ of field configurations
admits a cotangent complex generated by variations of the fields. Pairing these
variations through curvature and entropy flow produces a candidate closed
$3$--form, whose degree and symmetries align with $k=-1$.

In \cref{chap:plenum-stack}, we will:
\begin{enumerate}
  \item Construct $\Plenum$ as a derived Artin stack.
  \item Identify its cotangent complex $L_{\Plenum}$ explicitly in terms of
        variations of $(\PhiField,\vField,\SField)$.
  \item Construct a $(-1)$--shifted symplectic form $\omega_{\Plenum}$
        expressed locally as a derived $3$--form built from the lamphrodynamic
        fields.
  \item Verify that $\omega_{\Plenum}$ is closed and nondegenerate.
\end{enumerate}

This is a central technical achievement of Volume~I.

%--------------------------------------------------------------------
\section{Why \texorpdfstring{$k=-1$}{k = -1}?}
%--------------------------------------------------------------------

A natural question is: why specifically $k=-1$ for \rsvpabbr{}? Why not $k=0$
or $k=-2$?

\begin{itemize}
  \item The case $k=0$ would produce an \emph{even} symplectic form. This is
  appropriate for classical mechanics or certain topological sigma models, but
  not for gauge theories with ghost content.

  \item The BV formalism requires an odd symplectic form so that the BV bracket
  has the correct parity. This forces $k$ to be odd and negative.
\end{itemize}

Moreover, gauge symmetries and redundancy among lamphrodynamic fields indicate
that $\Plenum$ is not a naive moduli space but a derived stack with nontrivial
homotopy in negative degrees. The canonical choice compatible with BV is
$k=-1$, aligning with the degree of the ghost sector and with the reductions
arising in AKSZ.

\begin{remark}
In some exotic generalizations, higher negative shifts may appear (e.g.
topological field theories or higher spin models). For \rsvpabbr{}, the
entropic and lamphrodynamic structure strongly suggest $k=-1$ as the natural
choice.
\end{remark}

%--------------------------------------------------------------------
\section{Shifted Symplecticity and AKSZ}
%--------------------------------------------------------------------

The AKSZ construction, introduced in \cite{AKSZ}, produces solutions of the
classical master equation from a target derived stack equipped with a shifted
symplectic structure. The recipe is:
\begin{enumerate}
  \item Choose a source manifold $M$ (e.g.\ spacetime).
  \item Choose a target derived stack $F$ with an $n$--shifted symplectic
    structure.
  \item Form the mapping stack $\Map(M,F)$; by PTVV, this inherits an
    $(n-\dim M)$--shifted symplectic structure.
  \item For $n-\dim M = -1$, obtain an odd symplectic structure on
    $\Map(M,F)$, suitable for BV quantization.
\end{enumerate}

In \rsvpabbr{}, $F$ will ultimately be the lamphrodynamic target $\Plenum$ and
$M$ the spacetime manifold. The resulting $(-1)$--shifted symplectic structure
on $\Map(M,\Plenum)$ provides the geometric foundation for the BV action
$S_{\rsvpabbr}$.

%--------------------------------------------------------------------
\section{Summary and Forward Outlook}
%--------------------------------------------------------------------

Shifted symplectic geometry generalizes classical symplectic structures to a
homotopical context and provides the correct language for BV and AKSZ
quantization. The case $k=-1$ corresponds precisely to odd symplectic geometry
required for lamphrodynamic gauge theory.

In the next chapter, \cref{chap:aksz-bv}, we will review the AKSZ formalism in
detail, emphasizing how shifted symplectic structures on target stacks
determine BV actions, and how the constraint $n-\dim M=-1$ leads naturally to
lamphrodynamic dynamics on $\Plenum$.

